یہ کتاب دسویں جماعت کے طلبہ و طالبات  کے لئے  لکھی گئی ہے۔ یہ اُس کاوش کا حصہ ہے،  جو پاکستان میں سائنسی تعلیم  اردو اور دیگر علاقائی زبانوں میں مہیا کرنے کی خاطر کی جا رہی ہے۔ امید کی جاتی ہے کہ اس طرح کی کتب ہمارے نوجوان نسل کے لئے  تعلیم کے نیے راستے پید کرے  گی۔ 

یہ کتاب  کامسیٹ یونیورسٹی، اسلام آباد کے شعبہ طبیعیات  کے  درج ذیل طلبہ و طالبات نے  اسسٹنٹ پروفیسر زرقا  زاہد  کے زیر نگرانی     تشکیل دی۔طلبا  کو چار گروہوں میں تقسیم کیا گیا جن کی سربراہی    سید شاہمیر کازمی، مرزا اکبر اعوان، محمد عمار خان نیازی، اور عزا محمد کاظمی نے کی۔

Adiqa Hamza, Amber Riaz Khan, Ammara, Anila , Arfa Khatoon , Arwa Hassan, Asad, Asadullah, Asif, Atiqa Qarar, Awais, Ayesha Jamshaid, Fahad Khan, Hafiza Maryam, Haleema, Halima Ramzan., Hamza Hayat Khan, Hani Rizvi, Iqra Bashir , Jaweria Hosaini, Kainat, Khadija, Khadija Gul Sharif, Laiba Imran, Laraib Niaz, liza Sajjid, Lizza, Lutfullah, M Zakir, Maan Malik, Mahnoor Mehmood, Malik Hunza Awan, Misbah Batool, Misbah Shaheen, Muhammad Moosa, Muhammad Sohail, Nida Barkat, Noor ul Wara, Owais, Qurat ul Ain, Raza, Saadullah Pirzada, Sabeen, Sabeen Abid, Sana Noor, Sana Shabbir, Shan Tariq, Syed Zakir Shah, Tahoor, Wahaj, Yumna Khan, Zahra Fatima, Zahra Kazmi, Zain Abbas


زرقا زاہد\\
\اردوعدد{22}  جون \سن{2023}

